% =====================================================================
%  CV de Thomas FUZELLIER
% ---------------------------------------------------------------------
%  Ce fichier contient le CONTENU du CV : modifie-le librement.
%  Le DESIGN (couleurs, tailles, marges) est dans cv-style.sty.
%
%  Compilation : pdflatex cv_Thomas_Fuzellier.tex   (ou "Build LaTeX project" dans VS Code)
%
%  Rappels LaTeX :
%    - les caractères & % $ # _ s'écrivent \& \% \$ \# \_
%    - "--" donne un tiret de date (–)
%    - "~" est une espace insécable (ex. avant les deux-points)
%    - \textbf{...} met en gras (mots-clés des missions)
%    - une ligne commençant par % est un commentaire (ignorée)
% =====================================================================
\documentclass[10pt]{article}
\usepackage{cv-style}
\hypersetup{pdftitle={CV Thomas FUZELLIER}, pdfauthor={Thomas FUZELLIER}}

\begin{document}
\noindent
% #####################################################################
%  COLONNE GAUCHE (fond bleu)
%  Les rubriques s'étalent automatiquement sur toute la hauteur.
% #####################################################################
\begin{cvsidebar}

\cvphoto{photo.jpg}

% --- Contact : \sidecontact{icône}{texte} ------------------------------
\sidesection{Contact}
\sidecontact{\faEnvelope}{\href{mailto:thomas.fuzellier@outlook.fr}{thomas.fuzellier@outlook.fr}}
\sidecontact{\faPhone}{\href{tel:+33610484137}{+33~(0)~6~10~48~41~37}}
\sidecontact{\faGlobe}{\href{https://fuzellier-portfolio-website.vercel.app}{fuzellier-portfolio-website.vercel.app}}
\sidecontact{\faMapMarker}{Mobilité~: France}
\sidecontact{\faCar}{Permis B}
\sidecontact{\faFlag}{Nationalité française}

% --- Compétences : \sidegroup{Catégorie}{liste} --------------------------
%  Ordre : du plus important au moins important pour le poste visé.
\sidesection{Compétences}
\sidegroup{Langages}{C/C++, Python, Java, SQL}
\sidegroup{Systèmes embarqués}{STM32, Arduino, bus CAN}
\sidegroup{Électronique}{Conception PCB (KiCad), électronique embarquée}
\sidegroup{IA / Machine Learning}{LLM (GPT, LLaMA), NLP, Transformers, Prompt Engineering}
\sidegroup{Web \& API}{API REST, HTTP, HTML, CSS, JavaScript}
\sidegroup{Outils}{Git, Linux, Docker, Matlab, Fusion~360, PyTorch, HuggingFace, Pandas, NumPy}

% --- Qualités (rubrique optionnelle : enlève les % pour l'activer) -----
% \sidesection{Qualités}
% Qualité 1, Qualité 2, Qualité 3\par

% --- Langues : \sideitem{Langue}{Niveau} --------------------------------
\sidesection{Langues}
\sideitem{Anglais}{TOEIC 810}
\sideitem{Espagnol}{intermédiaire}

% --- Centres d'intérêt --------------------------------------------------
\sidesection{Centres d’intérêt}
\sidegroup{Menuiserie}{Fabrication de mobilier artisanal}
\sidegroup{Jeux}{Jeux de cartes, de plateau et vidéo}
\sidegroup{Sport}{Street Lifting, Tennis}

\end{cvsidebar}%
% #####################################################################
%  COLONNE DROITE (fond blanc)
% #####################################################################
\begin{cvmain}

% --- En-tête : \cvheader{Prénom}{NOM}{Intitulé} -------------------------
\cvheader{Thomas}{FUZELLIER}{Ingénieur Systèmes Embarqués}

% --- Accroche (supprimer la commande pour l'enlever) --------------------
\cvsummary{Ingénieur ENSEA en systèmes embarqués, avec plus d’un an 
d’expérience du firmware STM32 (C/C++) à la mise en service et au 
diagnostic sur le terrain. Profil polyvalent, du logiciel à l’électronique, 
avec une ouverture à l’IA appliquée.}

% =====================================================================
%  EXPÉRIENCES
%  \cvexperience{Titre}{Dates}{Entreprise}{Lieu}{Contexte en une ligne}
%  puis la liste des missions, et \cvsep entre deux expériences.
% =====================================================================
\cvsection{Expériences}
\begin{cvtimeline}

\cvexperience
  {Ingénieur Systèmes Embarqués -- Consultant}
  {Sept. 2026 -- aujourd’hui}
  {Ethic Opex - GCK}
  {Cournon, France}
  {Mission au service après-vente~: diagnostic, fiabilisation et amélioration continue des véhicules.}
\begin{cvitems}
  \item \textbf{Diagnostic hardware et software} des dysfonctionnements remontés par les utilisateurs
  \item \textbf{Solutions correctives}~: réalignement logiciel, \textbf{calibration}, instrumentation d’essai
  \item Analyse de données avec des \textbf{outils réseau et CAN}, \textbf{essais dynamiques} de validation
  \item \textbf{Rapports d’essais} et fiches d’intervention, remontée d’anomalies au \textbf{bureau d’études}
  \item Animation de \textbf{formations} destinées aux partenaires réparateurs
\end{cvitems}

\cvsep
\cvexperience
  {Ingénieur Systèmes Embarqués}
  {Janv. 2025 -- Janv. 2026}
  {Semi Lift}
  {Toulon, France}
  {PME du secteur des ascenseurs, internalisant son développement électronique.}
\begin{cvitems}
  \item Développement et maintenance du \textbf{firmware C/C++} de contrôle d’ascenseur sur \textbf{microcontrôleurs STM32}
  \item Développement de \textbf{drivers} pour l’interfaçage de capteurs et d’actionneurs
  \item \textbf{Intégration et validation} des capteurs sur systèmes d’ascenseur
  \item \textbf{Mise en service sur site} et dépannage des installations
  \item Développement d’\textbf{outils Python} pour automatiser les tests et analyser les données de fonctionnement
  % Mission retirée car peu utile pour un poste embarqué (enlever le % pour la remettre) :
  % \item Gestion et maintenance du site internet de l’entreprise
\end{cvitems}

\cvsep
\cvexperience
  {Stage Ingénieur NLP / LLM appliqué à la santé}
  {Mai -- Déc. 2024}
  {CTB - Laboratoire MEDAL}
  {Université de Madrid, Espagne}
  {Structuration automatique de notes cliniques confidentielles en espagnol
   à l’aide de LLM (Llama).}
\begin{cvitems}
  \item Conception d’un \textbf{pipeline NLP end-to-end} en Python~: extraction, normalisation et structuration des notes
  \item Intégration de \textbf{modèles Llama via API REST}~: prompt engineering, gestion du contexte et des tokens, optimisation et benchmark sur notes réelles
  \item Intégration d’un \textbf{knowledge graph} pour relier les informations médicales extraites
  \item \textbf{Industrialisation}~: architecture modulaire, tests unitaires, documentation technique
\end{cvitems}

\cvsep
\cvexperience
  {Stage Ingénieur Systèmes embarqués \& instrumentation}
  {Juin -- Août 2023}
  {Stratolia}
  {Palaiseau, France}
  {Nacelle d’observation de la Terre embarquée sous ballon stratosphérique.}
\begin{cvitems}
  \item Développement d’un \textbf{système de contrôle au sol} pour une caméra embarquée
  \item \textbf{Conception mécanique} et prototypage de la nacelle
  \item \textbf{Développement embarqué} sur microcontrôleur pour l’acquisition de données capteurs
  \item \textbf{Conception électronique} et réalisation de \textbf{PCB}
  \item Mise en place d’un \textbf{banc de test} simulant les contraintes du vol
\end{cvitems}

\end{cvtimeline}

% =====================================================================
%  FORMATION
%  \cveducation{Diplôme}{Dates}{Établissement}{Lieu}{Détails}
% =====================================================================
\cvsection{Formations}
\begin{cvtimeline}

\cveducation
  {Diplôme d’ingénieur}
  {Sept. 2020 -- Déc. 2024}
  {ENSEA — École Nationale Supérieure de l’Électronique et de ses Applications}
  {Cergy, France}
  {Ingénieur généraliste, option Informatique et Système}

\cvsep
\cveducation
  {Classe préparatoire ATS}
  {Sept. 2019 -- Juin 2020}
  {Lycée des Remparts}
  {Marseille, France}
  {Mathématiques et physique}

\cvsep
\cveducation
  {DUT GEII}
  {Sept. 2017 -- Juin 2019}
  {Université d’Aix-Marseille}
  {Marseille, France}
  {Génie Électrique et Informatique Industrielle}

\end{cvtimeline}

% =====================================================================
%  PROJETS
%  \cvproject{Titre}{Date}{Lieu (peut rester vide)}{Technologies}{Description}
% =====================================================================
\cvsection{Projets}
\begin{cvtimeline}

\cvproject
  {Générateur de candidatures assisté par IA}
  {Projet personnel}
  {\href{https://github.com/thomas-fuzellier/AI-Job-Application-Optimizer}{GitHub}}
  {Python · Tkinter · API Groq (LLM) · pytest}
  {Application générant CV adapté, lettre PDF et suivi Excel à partir d’une offre d’emploi.}
\begin{cvitems}
  \item \textbf{Pipeline LLM}~: détection de la langue, choix du CV et des paragraphes pré-rédigés
  \item IA limitée à la \textbf{sélection de contenus validés}, sans texte généré, pour un résultat maîtrisé
  \item \textbf{Architecture modulaire}, export \textbf{PDF} et \textbf{Excel}, \textbf{tests unitaires} avec API simulée
\end{cvitems}

\end{cvtimeline}

\end{cvmain}
\end{document}